\documentclass[10pt]{article}
\usepackage[utf8]{inputenc}
\usepackage{infocomp}
\usepackage{times}
\usepackage{amsmath}
\usepackage{amssymb}
\usepackage[T1]{fontenc}
\usepackage[brazilian]{babel}
\usepackage{graphicx}
\usepackage{hyperref}
\usepackage{subfigure}
\usepackage{enumerate}
\usepackage{caption}
\sloppy
\renewcommand{\captionfont}{\footnotesize}
\renewcommand{\captionlabelfont}{\footnotesize \bfseries}

\newtheorem{exemplo}{Exemplogfsdgdgdfsgdfsgsdfgdsfgdsfgdsfgdsf}

\address{UFLA - Universidade Federal de Lavras\\
        DCC - Departmento de Ci\^encia da Computa\c{c}\~ao\\
        P.O. Box 3037 - Campus da UFLA
		37200-000 - Lavras (MG)- Brazil\\
        $^1$\url{(first,second)@dcc.ufla.br}\\
        $^2$\url{third@comp.ufla.br}}

\title{Model for Articles to INFOCOMP}

\author{
        First Author of the Article$^1$ \\
        Second Author of the Article$^1$ \\
        Third Author of the Article$^2$
}

\abstract{This paper describes a template for INFOCOMP manuscripts.
INFOCOMP receives papers in continuous flow and publish all accepted
papers, however not necessarilly in the following issue.}

\keywords{template, manuscripts, model, INFOCOMP.}

\receivedate{July 19th, 2011}

\acceptdate{September 1st, 2011}

\begin{document}

\maketitle


\newpage

\section{Introduction}

In edition INFOCOMP Journal 2008 shall be published only papers
written in English. This change will allow that articles can be read
by several researchers from different countries. Thus, documents
submitted in Portuguese to be translated into good English.

The used principal source must be Times or Times
New Roman, size 10 points.

The space between lines must be of 12 points. After each paragraph
must be added an additional space of 3 points.

The used format of page is letter (to letter). The maximum amount of
pages is ten. The edges to be used are:
\begin{itemize}
\item top: 3cm;
\item botton, left and right: 2.5cm.
\end{itemize}

\section{First Page}

The first page must contain, in the following sequence:

\begin{itemize}
\item Title of the paper in source Helvetica or Arial, size 14 points
in boldface. Space before 12 points and after 24 points. Single
spacing between the lines;
\item Name of the authors in source Times or Times New Roman, size 10
points in high box, centered, space after of 3 points;
\item Address of the authors, in source Times or Times New Roman, size
10, centered, space after of 3 points;
\item Abstract of the paper in source Times or Times New Roman, size
10, with jib of 1cm to the left and the right, space after of 3 points
(the title ``Abstract:'' it must be in boldface);
\item Keywords keys of the paper, with identical formatting to the Abstract.
\end {itemize} 


\section{Sections}

Title of sections and subsection must be in source Arial or Helvetic,
size 10 points, in boldface. It must have an additional space of 9
points before the paragraph of heading and 3 after. The sections and
subsection must be numbered in the format: \\

1 Title of section 1

1.1 Title of Subsection 1.1

1.2 Title of Subsection 1.2

2 Title of section 2

2.1 Title of section 2.1\\

The first paragraph of each section does not have to be indent. The
too much paragraphs must be indent in 0.5cm. The text of the paper
must be in two separate columns for 0.5cm.

\section{Figures and tables}

Figures and tables can exceed the limits of the columns (but not them
of the edges). The heading of one given figure (to see Figure
\ref{fig:1}) must come below of the same one and centered.

\begin{figure}[!htbp]
  \begin{center}
    \includegraphics[scale=3]{exfigura}
    \caption{Figure of paper}
    \label{fig:1}
  \end{center}
\end{figure}

Before each figure must have a space of 6 points.  Between the figure
and the legend, a space of 3 points. After the legend must have a
space of 6 points. The legend must be in Times or Times New Roman,
size 9 points.

The table heading must come above of the same ones (to see Table
\ref{table1}).

\begin{table}[!htb]
  \begin{center}
    \caption{Example of Table}
    \label{table1}
    \vspace{0.2cm}
    \footnotesize
    \begin{tabular}{|c|c|c|c|c|c|}
      \hline
      $U$ & $vitA$ & $vitC$ & $vitD$ & $prot$ & $lip$ \\
      \hline
      \hline
      $d_1$ & 1 & 3 & 4 & 2 & 3\\
      $d_2$ & 1 & 3 & 3 & 3 & 2\\
      $d_3$ & 1 & 3 & 4 & 3 & 1\\
      $d_4$ & 3 & 5 & 2 & 5 & 2\\
      $d_5$ & 4 & 5 & 2 & 5 & 1\\
      $d_6$ & 3 & 5 & 2 & 3 & 4\\
      $d_7$ & 4 & 4 & 1 & 3 & 2\\
      \hline
    \end{tabular}
    %\vspace{-0.7cm}
  \end{center}
\end{table}

\section{Final version of the paper}

In final version of the paper, the authors must send the follows
files: 
\begin{enumerate}[a)]
\item Final paper in pdf;
\item Code source {\bf only in Latex};
\item The figures should sent separately, only in the formats: jpg,
png, eps and bmp.
\item All files must compressed within a file in formats: zip, tar,
rar or another.
\end{enumerate}

In the final version all the reviews must be considered to modify the
text. Papers should be submitted to INFOCOMP using the JEMS system of
Brazilian Computer Society in the page:\\
\url{https://submissoes.sbc.org.br/infocomp}

\section{Bibliograph References}

The bibliographical references will have to follow Norma NBR6023 of
the ABNT-2000, listed in sequence alphabetical, with obligator
numeration. 

The citations in the text must be made to leave of the numeration of
the bibliographical reference. The references must be inserted in file
bibtex, such as in the {\em example.bib} appended. 

 Example, ``In \cite{Correia}, is presented a taxonomy for medium
access control protocols in wireless sensor networks ...''.

\bibliography{example}
\end{document}
